%% MeTTaIL companion note for leanOSLF
%%
%% Build:
%%   pdflatex MeTTaIL.tex

\documentclass[11pt]{article}
\usepackage[T1]{fontenc}
\usepackage[utf8]{inputenc}
\usepackage{lmodern}
\usepackage{geometry}
\geometry{margin=1in}
\usepackage{hyperref}
\usepackage{longtable}
\usepackage{booktabs}
\usepackage{listings}
\usepackage{amssymb}
\lstset{
  basicstyle=\ttfamily\small,
  breaklines=true,
  frame=single,
  columns=fullflexible
}

\newcommand{\code}[1]{\texttt{#1}}
\newcommand{\codepath}[1]{\path{#1}}

\title{MeTTaIL + MeTTaCore in Lean and Rust\\
\large A Focused Companion to \code{leanOSLF.tex}}
\author{Mettapedia Project Notes (Draft)}
\date{\today}

\begin{document}
\maketitle

\begin{abstract}
MeTTaIL is a formally verified rewrite-rule language framework embedded in Lean~4.
It provides a syntax for pattern-matching rewrite rules with typed premises, a
declarative one-step semantics, a proved-sound-and-complete executable engine,
and an automatic derivation of an OSLF behavioral type system for any
\code{LanguageDef} instance.  This note covers the MeTTa-facing slice of the
formalization: the MeTTaIL and MeTTaCore Lean libraries, the MeTTaMinimal and
MeTTaFull instance endpoints, the Rust execution bridge, and a machine-checked
formalization of MQ-calculus (Stay \& Meredith 2026) as a concrete verification
target exercising the full stack.
\end{abstract}

\section{Current Coverage}

\begin{longtable}{lrr}
\toprule
Component & Lines & Sorries \\
\midrule
\endhead
\code{Mettapedia/OSLF/MeTTaIL/} & 5,250 & 0 \\
\code{Mettapedia/OSLF/MeTTaCore/} & 4,666 & 0 \\
\code{Mettapedia/Languages/MeTTa/HE/} & 2,774 & 0 \\
\code{Mettapedia/OSLF/PeTTa/} & 3,921 & 0 \\
\code{Mettapedia/Logic/GovernanceReasoning/} (MeTTa) & 1,140 & 0 \\
\code{Framework/MeTTaMinimalInstance.lean} & 700 & 0 \\
\code{Framework/MeTTaFullInstance.lean} & 346 & 0 \\
\midrule
Total (MeTTa-focused Lean slice) & 18,797 & 0 \\
\bottomrule
\end{longtable}

\section{Lean Side: What Is Formalized}

\subsection{MeTTaIL Syntax and Rule Language}

The syntax layer is the shared grammar from which all MeTTaIL language instances
are built: any \code{LanguageDef} is assembled from these types, so getting them
right is a prerequisite for correctness of every higher layer.
Core syntax and rule structures (locally nameless patterns, premises, equations,
rewrites, language definitions) are in
\codepath{Mettapedia/OSLF/MeTTaIL/Syntax.lean}.

\begin{lstlisting}[language={},caption={Core MeTTaIL structures (excerpt)}]
inductive CarrierKind where
  | ast | tokenLabel | tokenRaw | tokenProof
  | tokenPath | builtinInt | builtinString | builtinBool

structure TypeDecl where
  name : String
  carrier : CarrierKind := .ast

inductive Pattern where
  | bvar : Nat -> Pattern
  | fvar : String -> Pattern
  | apply : String -> List Pattern -> Pattern
  | lambda : Option String -> Pattern -> Pattern      -- authored binder name
  | multiLambda : Nat -> List String -> Pattern -> Pattern
  | subst : Pattern -> Pattern -> Pattern
  | collection : CollType -> List Pattern -> Option String -> Pattern

inductive SyntaxPatternOp where             -- Rust-style metasyntax
  | var : String -> SyntaxPatternOp
  | sep : String -> String -> Option SyntaxPatternOp -> SyntaxPatternOp
  | zip : String -> String -> SyntaxPatternOp
  | map : SyntaxPatternOp -> List String -> List SyntaxItem -> SyntaxPatternOp
  | opt : List SyntaxItem -> SyntaxPatternOp

inductive Premise where
  | freshness : FreshnessCondition -> Premise
  | congruence : Pattern -> Pattern -> Premise
  | relationQuery : String -> List Pattern -> Premise
  | forAll : String -> String -> Premise -> Premise    -- quantified premises

structure LanguageDef where
  name : String
  types : List TypeDecl                     -- carrier-aware type declarations
  terms : List GrammarRule
  equations : List Equation
  rewrites : List RewriteRule
  congruenceCollections : List CollType := []  -- opt-in per language
  logic : List LogicDecl := []
  oracles : List OracleDecl := []
\end{lstlisting}

Key design choices, each reaching council quorum:
\begin{itemize}
\item \textbf{Carrier-aware types} (\code{TypeDecl} with \code{CarrierKind}):
  Metamath surface syntax requires raw lexical carriers (\code{![raw] as Sym},
  \code{![label] as Label}). Rust already supports these via
  \code{LangType.carrier\_kind}; Lean now matches.
\item \textbf{Authored binder names} (\code{Pattern.lambda : Option String}):
  The human-facing name is preserved for export and diagnostics.
  The locally nameless core uses \code{none}; proofs are unaffected.
\item \textbf{Syntax operators} (\code{SyntaxPatternOp}):
  Rust's \code{*sep}, \code{*zip}, \code{*map}, \code{*opt} are first-class
  authored data, not flattened to terminal/nonTerminal strings.
\item \textbf{Quantified premises} (\code{Premise.forAll}):
  Required for scope-extrusion rules like
  \code{xs.*map(|x| x \# ...rest)}.
\item \textbf{Default congruence is empty}:
  Each language opts in explicitly; no hidden reduction behavior.
\item \textbf{Semantic validation} (\code{LanguageDef.validate}):
  Checks for unknown types, duplicate constructors, unbound syntax parameters,
  unknown relations, and arity mismatches---all at definition time.
\end{itemize}

\subsection{The \code{languageDef!} Authored Surface}

The Rust side uses a \code{language!} procedural macro to let humans author
language definitions in a single readable block.  Lean now has a matching
\code{languageDef!} macro that parses the \emph{same} surface syntax into the
\emph{same} \code{LanguageDef} type---no separate ``surface AST'' is introduced.
The macro \emph{is} the surface layer; the single \code{LanguageDef} is the
core.

\begin{lstlisting}[language={},caption={Lean \code{languageDef!} example
  (\codepath{LanguageDefDSLTests.lean})}]
def smokeLang : LanguageDef :=
  languageDef! {
    name : "SmokeLang"
    types {
      Expr
      Env
      Result
      ![raw] as Sym                       -- carrier-aware type
    }
    terms {
      SymLeaf . tok:Sym |- tok : Expr;
      Lam . ^ x . body : [Expr -> Expr]   -- binder name preserved
        |- "lam" body : Expr;
    }
    equations {
      EqSym . tok:Sym |- SymLeaf(tok) = SymLeaf(tok);
    }
    rewrites {
      Beta . env:Env | knownSymbol(SymLeaf(tok))
        |- Lam(SymLeaf(tok)) ~> SymLeaf(tok);
    }
    logic {
      relation knownSymbol(Expr);
      rule "knownSymbol(x) <-- true";
    }
    oracles {
      external evalExpr(Expr, Env) -> Result;
    }
    congruenceCollections { HashBag }
  }
\end{lstlisting}

The Rust and Lean term-declaration surfaces are intentionally identical:
%
\begin{center}
\begin{tabular}{@{}l@{\qquad}l@{}}
\textbf{Rust \code{language!}} & \textbf{Lean \code{languageDef!}} \\[2pt]
\code{PPar~.~ps:HashBag(Proc)} & \code{PPar~.~ps:HashBag(Proc)} \\
\code{~~|- "\{" ps.*sep("|") "\}" : Proc;}
  & \code{~~|- "\{" ps.*sep("|") "\}" : Proc;}
\end{tabular}
\end{center}

Three additional layers support the authored surface:
\begin{enumerate}
\item \textbf{Validation.}
  \code{LanguageDef.validate} checks types, constructors, relations,
  syntax parameters, and duplicates at definition time (kernel-checked by
  \code{native\_decide} in test examples).
\item \textbf{Fail-fast core bridge.}
  \code{CoreSyntaxBridge.specToCoreLanguage} returns \code{Except String}
  and rejects unsupported features (e.g.\ \code{forAll} premises,
  non-AST carriers, syntax metasyntax) rather than silently erasing them.
\item \textbf{Faithful export.}
  \code{Export.renderLanguageWithUserSyntax} preserves authored binder names,
  \code{...rest} collection remainders, and syntax operators---no invented
  \code{C\_} prefixes or synthetic binder names.
\end{enumerate}

\noindent
The overall architecture is a single pipeline with two host languages:

\begin{center}
\begin{tabular}{c}
\textbf{Human-Authored Source} \\
Rust \code{language!} / Lean \code{languageDef!} \\[4pt]
$\downarrow$ \\[4pt]
\textbf{ONE \code{LanguageDef}} \\
(types, terms, equations, rewrites, logic, oracles) \\[4pt]
$\swarrow$ \hspace{4em} $\searrow$ \\[4pt]
\begin{tabular}{c}
\textbf{Rust Runtime} \\
MORK/MM2, Ascent, \\
tree-sitter parser
\end{tabular}
\hspace{2em}
\begin{tabular}{c}
\textbf{Lean Verification} \\
Proofs, OSLF typing, \\
validated export
\end{tabular}
\end{tabular}
\end{center}

\subsection{RhoCalc as the \code{LanguageDef} Parity Pressure Test}

The $\rho$-calculus (Meredith \& Stay) is the flagship language instance for both
the Rust \code{language!} macro and the Lean \code{languageDef!} surface.  The
Rust definition in \codepath{hyperon/mettail-rust/languages/src/rhocalc.rs}
(${\sim}380$ lines) covers the full process calculus including native arithmetic,
string operations, and fold terms.

Lean now has a directly authored process fragment
(\codepath{Languages/ProcessCalculi/RhoCalculus/LanguageDefDSL.lean})
covering the pure process kernel:

\begin{lstlisting}[language={},caption={RhoCalc process fragment in Lean
  \code{languageDef!} form (excerpt)}]
def rhoCalcProcessFragment : LanguageDef :=
  languageDef! {
    name : "RhoCalc"
    types { Proc; Name }
    terms {
      PZero  . |- "{}" : Proc;
      PDrop  . n:Name |- "*" "(" n ")" : Proc;
      PPar   . ps:HashBag(Proc) |- "{" ps.*sep("|") "}" : Proc;
      POutput. n:Name, q:Proc |- n "!" "(" q ")" : Proc;
      PInputs. ns:Vec(Name), ^[xs].p:[Name* -> Proc]
        |- "(" *zip(ns,xs).*map(|n,x| n "?" x).*sep(",") ")"
           "." "{" p "}" : Proc;
      NQuote . p:Proc |- "@" "(" p ")" : Name;
      PNew   . ^[xs].p:[Name* -> Proc]
        |- "new" "(" xs.*sep(",") ")" "in" "{" p "}" : Proc;
    }
    equations {
      QuoteDrop . N:Name |- NQuote(PDrop(N)) = N;
    }
    rewrites {
      Exec    . P:Proc |- PDrop(NQuote(P)) ~> P;
      ParCong . | S ~> T
        |- PPar({S, ...rest}) ~> PPar({T, ...rest});
      NewCong . | S ~> T
        |- PNew(^[xs].S) ~> PNew(^[xs].T);
    }
    congruenceCollections { HashBag }
  }
\end{lstlisting}

All structural properties are kernel-checked: \code{LanguageDef.validate}
returns \code{[]} (proved by \code{native\_decide}), and export substring
checks confirm that the Lean-rendered text matches the Rust surface.

\paragraph{Honest boundary.}
Two Rust rules cannot yet be authored in Lean:
\begin{itemize}
\item \textbf{Comm} requires \code{*zip(ns,qs).*map(...)} in \emph{pattern}
  position (not just syntax-pattern position) and a direct \code{(eval cont ...)}
  substitution form.
\item \textbf{Extrude} requires \code{xs.*map(|x| x \# ...rest)} in
  \emph{premise} position---quantified freshness over a collection.
\end{itemize}
These are the next DSL extensions needed for full parity.

\subsection{Premise-Aware Declarative Semantics and Engine Equivalence}

The biconditional between the declarative relation and the executable engine
is the central correctness guarantee of MeTTaIL: it means the engine can be
used for computation while the declarative relation is used for reasoning,
and the two are interchangeable by theorem.
The declarative one-step relation with relation-environment premises is
\code{DeclReducesWithPremises} in
\codepath{Mettapedia/OSLF/MeTTaIL/DeclReducesWithPremises.lean}, together with
soundness/completeness against the executable engine.

\begin{lstlisting}[language={},caption={Declarative semantics and engine correspondence}]
inductive DeclReducesWithPremises (relEnv : RelationEnv) (lang : LanguageDef) :
    Pattern -> Pattern -> Prop where
  | topRule : ...
  | congElem : ...

theorem engineWithPremisesUsing_sound :
    q in rewriteWithContextWithPremisesUsing relEnv lang p ->
    DeclReducesWithPremises relEnv lang p q

theorem engineWithPremisesUsing_complete :
    DeclReducesWithPremises relEnv lang p q ->
    q in rewriteWithContextWithPremisesUsing relEnv lang p

theorem declReducesWithPremises_iff_langReducesWithPremisesUsing :
    DeclReducesWithPremises relEnv lang p q <->
    q in rewriteWithContextWithPremisesUsing relEnv lang p
\end{lstlisting}

\subsection{OSLF Synthesis Hooks Used by MeTTa Instances}

Every \code{LanguageDef} automatically acquires a behavioral type system via
\code{langOSLF}: the modal Galois connection $\Diamond \dashv \Box$ and the OSLF
type system are derived, not hand-coded, which means language designers get
behavioral typing for free.
Language-level reduction, engine equivalence, Galois connection, and synthesized
OSLF type system are in
\codepath{Mettapedia/OSLF/Framework/TypeSynthesis.lean}.

\begin{lstlisting}[language={},caption={TypeSynthesis bridge theorems}]
theorem langReducesUsing_iff_execUsing (relEnv : RelationEnv) (lang : LanguageDef)
    (p q : Pattern) :
    langReducesUsing relEnv lang p q <-> langReducesExecUsing relEnv lang p q

theorem langGaloisUsing (relEnv : RelationEnv) (lang : LanguageDef) :
    GaloisConnection (langDiamondUsing relEnv lang) (langBoxUsing relEnv lang)

def langOSLF (lang : LanguageDef) (procSort : String := "Proc") :
    OSLFTypeSystem (langRewriteSystem lang procSort)
\end{lstlisting}

\subsection{MeTTaCore State Machine and Properties}

MeTTaCore models the MeTTa runtime as an explicit state machine over
multisets of atoms, giving precise machine-checked statements of progress
and determinism for the grounded-operation layer.
State and transition components are in
\codepath{Mettapedia/OSLF/MeTTaCore/State.lean} and
\codepath{Mettapedia/OSLF/MeTTaCore/RewriteRules.lean}; key properties are in
\codepath{Mettapedia/OSLF/MeTTaCore/Properties.lean}.

\begin{lstlisting}[language={},caption={Representative MeTTaCore properties}]
structure MeTTaState where
  input : Multiset Atom
  knowledge : Atomspace
  workspace : Multiset Atom
  output : Multiset Atom

theorem progress (s : MeTTaState) (a : Atom) (_ : a in s.workspace) :
    s.knowledge.insensitive a = true \/
    (s.knowledge.queryEquations a).card > 0

theorem grounded_confluence (op : String) (args : List Atom) :
    forall r1 r2,
    executeGroundedOp (.symbol op) args = some r1 ->
    executeGroundedOp (.symbol op) args = some r2 ->
    r1 = r2
\end{lstlisting}

\subsection{HE MeTTa: Computable Recursive Interpreter}

The Hyperon Experimental (HE) MeTTa semantics are formalized as a set of
six mutually recursive fuel-bounded functions in
\codepath{Mettapedia/Languages/MeTTa/HE/}, directly corresponding to
\code{metta}, \code{interpret\_expression}, \code{interpret\_atom},
\code{interpret\_symbol}, \code{interpret\_type}, and
\code{chain\_results} from the
\href{https://trueagi-io.github.io/hyperon-experimental/metta/}{HE specification}.

\begin{lstlisting}[language={},caption={HE MeTTa interpreter core (excerpt from Interpreter.lean)}]
mutual
  def mettaCall (fuel : Nat) (space : Space) (bindings : Bindings)
      (atom : Atom) (typ : Atom) : ResultSet
  def interpretExpression (fuel : Nat) (space : Space) (bindings : Bindings)
      (atom : Atom) (typ : Atom) : ResultSet
  def interpretAtom (fuel : Nat) (space : Space) (bindings : Bindings)
      (atom : Atom) (typ : Atom) : ResultSet
  -- ... (6 mutually recursive functions)
end
\end{lstlisting}

Supporting modules formalize the full HE runtime:
\begin{itemize}
\item \code{Types.lean}: ErrorCode, Bindings, ResultSet
\item \code{Space.lean}: Space, queryEquations, GroundedDispatch
\item \code{Matching.lean}: matchAtoms, mergeBindings
\item \code{TypeCheck.lean}: typeCast
\end{itemize}

\paragraph{Conformance.}
37 kernel-checked conformance theorems in \code{Conformance.lean} verify
concrete evaluations against expected results, all closed by \code{rfl}
(no axioms, no \code{native\_decide}).

\paragraph{OSLF pipeline export.}
The HE interpreter is also exported as an OSLF \code{LanguageDef}
(\code{HELanguageDef.lean}) and a \code{PremiseProgram}
(\code{HEPremises.lean}), enabling automatic generation of the Rust
Ascent backend.  The \code{PremiseProgram} specifies 19 premise
relations and their datalog rules as first-order data that the Lean
exporter renders to Ascent syntax.

\subsection{PeTTa: Logic-Programming MeTTa Semantics}

PeTTa formalizes MeTTa evaluation through a Prolog-inspired
logic-programming lens.  The core is \code{MeTTaEval}, an inductive
big-step evaluation relation in
\codepath{Mettapedia/OSLF/PeTTa/MeTTaEval.lean}:

\begin{lstlisting}[language={},caption={MeTTaEval inductive relation (excerpt)}]
inductive MeTTaEval : PeTTaSpace -> Pattern -> Atom -> List Pattern -> Prop where
  | symbolLookup (sp : PeTTaSpace) (sym : String) (results : List Pattern) :
      sp.queryEquations (.symbol sym) = results ->
      results != [] ->
      MeTTaEval sp (.symbol sym) undefinedType results
  | groundedDispatch ...
  | expressionInterpretation ...
\end{lstlisting}

This relation is the shared core that both HE and PeTTa agree on at the
answer level: the refinement proofs (Section~\ref{sec:governance}) show
that PeTTa refines HE evaluation.

Additional PeTTa modules include:
\begin{itemize}
\item \code{LPSoundness.lean}: LP-engine soundness for rewrite-only evaluation
\item \code{GroundedOracle.lean}: typed grounded-operation interface
\item \code{TypedEval.lean} / \code{TypeSystem.lean}: type-enriched evaluation
\item \code{Effects.lean} / \code{MinimalInstructions.lean}: effect tracking and instruction set
\item \code{TranslateExpr.lean}: expression-to-pattern translation
\end{itemize}

\subsection{Metaprogramming and Dialect Translation}
\label{sec:metaprog}

An experimental HE$\leftrightarrow$PeTTa translator
(\codepath{hyperon/translators/}) revealed a fundamental difference in
how the three MeTTa runtimes handle code-as-data during structural
rewriting.  All three runtimes can load source atoms into a space and
use \code{match} to inspect them without evaluation.  However, when a
match template \emph{reconstructs} the matched term with rewritten
subexpressions, the three runtimes diverge:

\begin{center}
\begin{tabular}{lll}
\toprule
Runtime & Output of chain$\to$let rewrite & Variable handling \\
\midrule
CeTTa (C) & \code{(= (sq \$x) (* (+ \$x 1) (+ \$x 1)))} & Unified/inlined \\
HE (Rust) & \code{(= (sq \$x) (* (+ \$x 1) (+ \$x 1)))} & Unified/inlined \\
PeTTa (Prolog) & \code{(= (sq \$x) (let \$y (+ \$x 1) (* \$y \$y)))} & Preserved \\
\bottomrule
\end{tabular}
\end{center}

All three outputs are \textbf{semantically equivalent} (they compute the
same function), but PeTTa produces the \textbf{syntactically faithful}
translation while HE/CeTTa produce a \textbf{beta-reduced} expansion.
The cause: MeTTa's \code{match} performs full unification across the
entire pattern.  When a template variable \code{\$v} binds to a source
variable \code{\$y}, all occurrences of \code{\$y} in the template are
also unified, effectively inlining the binding.  PeTTa's Prolog-based
matching uses \code{copy\_term} and fresh variable generation, which
preserves syntactic structure.

\paragraph{Implications.}
\begin{itemize}
\item Pattern mining, backward chaining, and other metaprogramming tasks
  that require structural inspection of code are naturally easier in
  PeTTa.  This may explain why PeTTa's pattern miner finds all results
  where HE's has edge-case failures.
\item For exact source-to-source translation, the translator should be
  written in PeTTa (primary) with a semantic-equivalence version in
  HE/CeTTa.
\item A future MeTTa-Pure could address this with an explicit quoting
  mechanism, de~Bruijn indices, or meta-level separation (as formalized
  in the $\rho$-calculus layer of MeTTaIL).
\end{itemize}

\paragraph{Bidirectional Prolog translator.}

The practical translator is implemented in SWI-Prolog
(\codepath{hyperon/translators/he\_to\_petta.pl},
\codepath{petta\_to\_he.pl}), where term manipulation is native and
variable scoping is handled correctly by Prolog's \code{copy\_term}.
Each translation rule is a single, auditable \code{translate\_term/2}
clause:

\begin{center}
\begin{tabular}{lll}
\toprule
Direction & Rule & Clause \\
\midrule
HE$\to$PeTTa & \code{chain}$\to$\code{let} & \code{translate\_term([chain,E,V,B], [let,V,TE,TB])} \\
HE$\to$PeTTa & \code{collapse-bind}$\to$\code{collapse} & \code{translate\_term([collapse-bind,X], [collapse,TX])} \\
HE$\to$PeTTa & \code{superpose-bind}$\to$\code{superpose} & \code{translate\_term([superpose-bind,X], [superpose,TX])} \\
HE$\to$PeTTa & \code{switch}$\to$\code{case} & \code{translate\_term([switch,S,B], [case,TS,B])} \\
HE$\to$PeTTa & \code{atom-subst}$\to$\code{let} & \code{translate\_term([atom-subst,A,V,T], [let,V,TA,TT])} \\
HE$\to$PeTTa & \code{function/return} & unwrap \\
HE$\to$PeTTa & \code{nop}$\to$\code{let \$\_} & side-effect wrapper \\
\midrule
PeTTa$\to$HE & \code{progn}$\to$nested \code{let \$\_} & 2- and 3-argument forms \\
PeTTa$\to$HE & \code{prog1}$\to$\code{let+let} & save-first pattern \\
PeTTa$\to$HE & \code{@<}$\to$\code{<s} & string comparison \\
\bottomrule
\end{tabular}
\end{center}

All rules recursively translate subexpressions via \code{maplist(translate\_term)};
shared constructs (\code{if}, \code{let}, \code{case}, \code{match},
\code{superpose}, \code{collapse}, arithmetic) pass through unchanged.

The translator is validated on the full corpus of real \code{.metta} files:

\begin{center}
\begin{tabular}{lrr}
\toprule
Direction & Files tested & Passed \\
\midrule
HE$\to$PeTTa (CeTTa test suite) & 93 & 93 \\
PeTTa$\to$HE (PeTTa examples + miner) & 344 & 344 \\
\midrule
\textbf{Total} & \textbf{437} & \textbf{437} \\
\bottomrule
\end{tabular}
\end{center}

A correct-by-construction MeTTa S-expression parser
(\codepath{hyperon/translators/metta\_parser.pl}) handles all MeTTa token
types: \code{\$variables}, \code{\%types\%}, \code{"strings"}, numbers,
operators, and nested S-expressions.  The lexer rule is minimal: tokens
are delimited by whitespace and parentheses only.

\paragraph{Towards a pure MeTTa translator.}

The Prolog translator works but lives outside MeTTa.  A pure MeTTa
implementation would close the bootstrapping loop (MeTTa translating
MeTTa).  The challenge documented in
\code{metaprogramming\_challenges.metta} is that MeTTa's
unification-based \code{match} prevents exact syntactic rewriting.
PeTTa's Prolog backend naturally avoids this via \code{copy\_term}, so a
PeTTa-native translator using \code{import\_prolog\_functions\_from\_file}
is the most promising path.  This would let PeTTa call the Prolog
translator rules directly from MeTTa code, combining MeTTa's space-based
program loading with Prolog's faithful term manipulation.

The translator corpus is at \codepath{hyperon/translators/}, with
\code{metaprogramming\_challenges.metta} documenting five approaches
and their failure modes.

\subsection{Governance Refinement Proofs}
\label{sec:governance}

The governance reasoning layer
(\codepath{Mettapedia/Logic/GovernanceReasoning/}) connects MeTTa
evaluation to deontic temporal logic via DDL+.  This is the
end-to-end chain: evaluate MeTTa code $\to$ decode results $\to$
interpret as deontic modalities $\to$ apply DTS theorems.

\paragraph{Shared interface.}
\code{LetStarInterface.lean} defines a \code{MeTTaLike} typeclass
parametric over any evaluator supporting rewrite-rule application.
Both HE and PeTTa instantiate it, so let* unfolding theorems are
proved once and work for both implementations.

\paragraph{HE refinement.}
\code{HERefinement.lean} bridges \code{MeTTaEval} to the governance
DTS obligations with explicit type/binding tracking. Seven semantic
refinement theorems are proved, including let* integration
(\code{he\_letStar\_full\_2}, \code{he\_letStar\_full\_3}).

\paragraph{PeTTa refinement.}
\code{PeTTaRefinement.lean} establishes the PeTTa-side DTS refinement
and let* unfolding at the PeTTa level.

\paragraph{HE canonical conformance.}
\code{HECanonicalConformance.lean} connects the canonical computable
HE interpreter to governance DTS theorems.  Seven conformance theorems
prove end-to-end: HE eval produces correct OB$\Rightarrow$PE results for
concrete governance scenarios (soaMoor, payRent, giveNotice), all closed
by \code{rfl} or \code{simp}.

All files: 0 sorries, 0 custom axioms.

\subsection{MeTTaMinimal and MeTTaFull Instance Endpoints}

These are the integration points where abstract MeTTaIL theorems become
concrete MeTTa-language guarantees: a developer who adds a new MeTTa feature
can check its behavioral soundness by extending one of these instances.
MeTTa instances export checker-soundness theorems with instance-specific atom
interpretations:

\begin{lstlisting}[language={},caption={MeTTa instance API endpoints}]
def mettaMinimalOSLF := langOSLF mettaMinimal "State"
theorem mettaMinimal_checkLangUsing_sat_sound_specAtoms :
    checkLangUsing ... = .sat -> sem ...

def mettaFull : LanguageDef := ...
def mettaFullOSLF := langOSLF mettaFull "State"
def mettaFullRelEnv : RelationEnv := ...
theorem mettaFull_checkLangUsing_sat_sound_specAtoms :
    checkLangUsing ... = .sat -> sem ...
\end{lstlisting}

\section{Type System Limitation: No Subject Reduction}
\label{sec:sr}

\subsection{Subject Reduction in Type Theory}

A type system satisfies \emph{subject reduction} (also called \emph{type preservation})
if typing is preserved under evaluation: whenever $a : T$ and $T$ reduces to $T'$,
we have $a : T'$.  Subject reduction is a foundational theorem in well-behaved dependent
type theories---it holds for Martin-L\"{o}f Type Theory, the Calculus of Constructions,
and Lean~4 itself.

\subsection{MeTTa's Typing Judgment}

MeTTa's type system, formalized in
\codepath{Mettapedia/OSLF/MeTTaCore/Types.lean},
uses a \emph{purely lookup-based} judgment:

\begin{lstlisting}[language={},caption={HasType judgment (excerpt from Types.lean)}]
inductive HasType (space : Atomspace) : Atom -> Atom -> Prop where
  | intrinsicSymbol (s : String) :
      HasType space (.symbol s) (.symbol "Symbol")
  | annotated (a ty : Atom) :
      typeAnnotation a ty in space.atoms ->
      HasType space a ty
  ...

def typeAnnotation (a ty : Atom) : Atom :=
  .expression [.symbol ":", a, ty]
\end{lstlisting}

The critical rule is \code{annotated}: the only way a term $a$ acquires a
non-intrinsic type $T$ is for the literal annotation \code{(: a T)} to be present
in the atomspace.  There is no inference rule that propagates types through
evaluation of the type expression itself.

\subsection{The Theorem}

\codepath{Mettapedia/OSLF/MeTTaCore/SubjectReduction.lean} establishes:

\begin{lstlisting}[language={},caption={Subject reduction definitions and main theorem}]
def SubjectReduction
    (HasTy : Atomspace -> Atom -> Atom -> Prop)
    (Reduces : Atomspace -> Atom -> Atom -> Prop) : Prop :=
  forall (space : Atomspace) (a T T' : Atom),
    HasTy space a T -> Reduces space T T' -> HasTy space a T'

-- One-step syntactic rewrite: (= a b) in space.atoms
def AtomReduces (space : Atomspace) (a b : Atom) : Prop :=
  .expression [.symbol "=", a, b] in space.atoms

theorem metta_not_subject_reduction :
    not (SubjectReduction HasType AtomReduces)
\end{lstlisting}

The proof is by explicit counterexample. Define the atomspace
\[
  \mathcal{S} \;=\; \{\;\texttt{(: t (Vec n))},\;\; \texttt{(= (Vec n) (Vec 0))}\;\}.
\]
Three lemmas, each closed by kernel-checked computation (\code{decide}):
\begin{enumerate}
\item $t : \texttt{(Vec n)}$ in $\mathcal{S}$, via \code{HasType.annotated}.
\item $\texttt{(Vec n)}$ reduces to $\texttt{(Vec 0)}$ in $\mathcal{S}$,
  since \texttt{(= (Vec n) (Vec 0))} $\in \mathcal{S}$.
\item $t \not: \texttt{(Vec 0)}$ in $\mathcal{S}$: the only surviving case of the
  \code{HasType} induction is \code{annotated}, which would require
  \texttt{(: t (Vec 0))} $\in \mathcal{S}$---but the space contains no such atom.
  Lean~4 closes all other constructor cases automatically via index discrimination.
\end{enumerate}

The file also establishes the companion result \code{counterSpace\_not\_eval\_closed}:
the counterexample space fails the \emph{eval-closure} property---the atomspace-level
condition that \emph{would} suffice for subject reduction, namely that whenever
$\texttt{(: a T)} \in \mathcal{S}$ and $T$ reduces to $T'$, also
$\texttt{(: a T')} \in \mathcal{S}$.

\subsection{Significance}

This result is not a bug in MeTTa; it is a design consequence of separating
\emph{knowledge} (rewrite rules and annotations, stored in the atomspace) from
\emph{computation} (evaluation via pattern matching).  In MeTTa, types and terms are
both atoms evaluated by the same engine.  The engine can reduce a type expression, but
the typing relation does not track this: $\texttt{HasType}$ reflects the atomspace
at a fixed moment, not under its closure.

This contrasts with MLTT/CoC, where subject reduction is proved by induction on the
typing derivation: every rule carries a proof term that is transformed, not just
looked up, when the type reduces.  MeTTa's expressiveness (Turing-complete rewriting
over a uniform atom language) comes at the cost of this metatheoretic guarantee.
Eval-closure of the atomspace is a sufficient---though not necessary---runtime
condition for subject reduction to hold in a given MeTTa program.

\section{Rust Bridge Status (\code{hyperon/mettail-rust})}

The Rust side implements three MeTTa language backends, all registered
in the REPL and augmented with the meta-introspection and Python oracle
systems.

\subsection{Language Backends}

\begin{itemize}
\item \textbf{MeTTaHE} (780 lines):
  \codepath{languages/src/mettahe\_from\_lean.rs}---the HE MeTTa backend,
  generated from the Lean OSLF pipeline (\code{HELanguageDef.lean} +
  \code{HEPremises.lean}).  Defines 44 term constructors across 4 types
  (State, Instr, Atom, Space), 17 rewrite rules, and 19 premise relations
  implemented as Ascent datalog rules.  Eight builtin helper functions
  provide structural dispatchers for grounded arithmetic, equation lookup,
  type annotation search, and function-type resolution.

\item \textbf{MeTTaFullState} (1,618 lines, legacy):
  \codepath{languages/src/mettafull\_legacy.rs}---the original MeTTaFull
  backend with hand-written premise rules and \code{eval\_ground\_call\_core}
  recursive evaluator.  Retained for regression testing; the HE backend
  supersedes it for HE-conformant evaluation.

\item \textbf{PeTTa} (226 lines):
  \codepath{languages/src/petta\_from\_lean.rs}---currently delegates to
  MeTTaFullState with library aliases.  Planned for standalone
  reimplementation from Lean PeTTa specs.
\end{itemize}

All three are registered in \codepath{repl/src/registry.rs} with the
\code{OracleAugmentedLanguage} wrapper, which adds the meta-introspection
oracle (types/terms/rewrites/relations counts and membership queries) and
the gated Python oracle.

\subsection{Premise Program IR}

The MeTTaHE backend demonstrates the \code{PremiseProgram} pipeline:
premise semantics are specified as first-order datalog rules in Lean,
then rendered to Ascent syntax for the Rust backend.  Negation over
derived relations (forbidden by Ascent stratification) is handled via
positive complement relations (e.g., \code{typeNotMatchesMetaOrAtom}
instead of \code{!typeMatchesMetaOrAtom}).

\paragraph{Macro hygiene note.}
The \code{language!} proc macro generates parser functions with internal
variables (\code{pos}, \code{lhs}, \code{rhs}, \code{value}, etc.).
Term constructor field names that collide with these produce confusing
type errors.  The HE backend documents the required renames in its file
header.

\subsection{Current Passing Rust Tests}

\begin{lstlisting}[language={},caption={MeTTa Rust canaries (all passing)}]
relation_query_can_bind_fresh_output_vars
typecheck_uses_space_typing_relation
grounded_call_boolean_not_reduces
return_step_reaches_done_state
default_registry_contains_mettafull
\end{lstlisting}

\subsection{Metamath as a Three-Way Conformance Target}

Metamath is the first language where three independently developed artifacts
can be connected through \code{LanguageDef}:

\begin{center}
\begin{tabular}{c}
\code{metamath.rs} \\
(Rust \code{language!} macro) \\[4pt]
$\swarrow$ \hspace{6em} $\searrow$ \\[4pt]
\begin{tabular}{c}
\code{LanguageDefDSL.lean} \\
(Lean \code{languageDef!})
\end{tabular}
\hspace{2em}
\begin{tabular}{c}
\code{mm-lean4} \\
(verified implementation)
\end{tabular} \\[4pt]
$\searrow$ \hspace{6em} $\swarrow$ \\[4pt]
\textbf{Conformance theorem} \\
(future work)
\end{tabular}
\end{center}

The Rust \code{metamath.rs} defines the full Metamath surface---\code{Database},
\code{Stmt}, \code{MathString}, \code{ProofString}---with carrier-aware
type declarations (\code{![label] as Label}, \code{![raw] as Sym},
\code{![proofTok] as ProofTok}, \code{![path] as IncludePath}) and explicit
term constructors for every \code{\$c}, \code{\$v}, \code{\$f}, \code{\$e},
\code{\$a}, \code{\$p} declaration form.

The Lean side currently has a core-step fragment in
\codepath{Languages/Metamath/LanguageDSL.lean} exercising carrier types and
the direct authored DSL surface, but not the full Metamath grammar.

The \code{mm-lean4} project (\codepath{hyperon/metamath/mm-lean4/}) provides
a complete, sorry-free, axiom-free verification of the Metamath proof checker:
\begin{itemize}
\item \code{DeclarativeSpec.lean}: \code{CN} (constants), \code{VR} (variables),
  \code{Expr} (expressions), \code{Provable} (proof validity).
\item End-to-end biconditional:
  \code{verify\_parser\_acceptance\_iff\_spec\_provable}---raw bytes
  $\leftrightarrow$ \code{Spec.Provable}, with 0 sorries and 0 axioms.
\end{itemize}

The types correspond across all three artifacts:

\begin{center}
\begin{tabular}{lll}
\toprule
\textbf{Rust \code{metamath.rs}} & \textbf{Lean \code{LanguageDef}} & \textbf{mm-lean4 \code{DeclarativeSpec}} \\
\midrule
\code{Sym} (carrier: raw)   & \code{TypeDecl "Sym" .tokenRaw}   & \code{CN} / \code{VR} \\
\code{Label} (carrier: label) & \code{TypeDecl "Label" .tokenLabel} & labels in frame \\
\code{MathString}            & grammar rule chain                 & \code{Expr} \\
\code{Database}              & \code{DbMore} / \code{DbOne}       & parsed DB structure \\
\code{ProofString}           & proof token chain                  & proof witness \\
\bottomrule
\end{tabular}
\end{center}

\paragraph{What remains.}
The formal conformance theorem---showing that the \code{LanguageDef} rewrite
semantics agree with \code{mm-lean4}'s operational verification---is future
work.  The type correspondence above provides the skeleton; the hard part
is proving that the stack-machine semantics of Metamath proof verification
agree with the term-rewriting semantics of the \code{LanguageDef} engine.

\section{MQ-Calculus as a Formalized Verification Target}

MQ-calculus (Stay \& Meredith 2026) is a process calculus in which communication
events are measurement events: a synchronization on wire $i$ is not just a
handshake but a quantum measurement, selecting a continuation branch according
to Born-rule probabilities.  The grammar is deliberately minimal---six
constructors covering nil, parallel composition, fresh-wire allocation, gate
application, output, and measurement-branching input---which makes the full
formalization tractable while still exercising the semantically interesting cases.

MQ-calculus is included in this formalization as a concrete stress test for the
MeTTaIL stack.  It is not a special case; it is an instance of \code{LanguageDef}
and inherits the full MeTTaIL pipeline automatically: the syntax is given by the
\code{Process} inductive, the reduction relation \code{Reduces} is derived from
\code{CommReduction} and structural congruence \code{SC}, the denotational
semantics \code{denoteWith} is backend-parametric over \code{MQSemanticsBackend},
and the MORK cross-calculus coherence theorem
(\code{paper\_mork\_mq\_branching\_equiv}) connects MQ COMM branching to MORK
binary fold branching at the theorem level.

Three results distinguish the Lean formalization from a pencil-and-paper treatment.
First, the Born-rule normalization theorem \code{born\_prob\_sum\_one} is a
machine-checked proof from Mathlib's \code{PiLp.norm\_sq\_eq\_of\_L2}, not an
axiom.  Second, the denotational semantics is parameterized by an abstract
\code{MQSemanticsBackend}, so the same six denotational clauses are provably correct
for any compliant backend without re-proving them.  Third, every clause of
\S{}6.3 of the paper has a named Lean theorem in \code{PaperMap.lean} (see
Appendix~\ref{app:mq-lean}), giving an explicit audit surface that connects
narrative paper text to kernel-checked declarations.

The MQ formalization also provides realistic workload profiles for
\code{hyperon/mettail-rust}: COMM-heavy processes exercise nondeterministic
branching and collection-heavy parallel composition, which are the known Ascent
scaling pressure points documented in the engine performance notes.

\section{MeTTa-IL as a 2-Mode Adjoint Doctrine (MaTT)}
\label{sec:matt}

The OSLF framework establishes that MeTTa-IL is an instance of a
\emph{2-mode adjoint doctrine fragment} in the sense of
fibrational/doctrinal logic (Lawvere hyperdoctrines, Jacobs fibred
category theory).  This is the formal content of ``MeTTa-IL is MaTT.''

\paragraph{Theoretical position.}
MeTTa-IL-via-OSLF is not the syntactic multimodal type theory
of~\cite{gratzer2020multimodal} (which would require modal type formers
$\square_\mu A$, locks, and context extension).  It is instead the
\emph{categorical/doctrinal} flavour: a mode category over a base of
operational languages, equipped with a contravariant indexed predicate
functor and a modal adjunction $\lozenge \dashv \square$ at each
behavioral mode.  This places it squarely in the tradition of Lawvere
hyperdoctrines and fibrational semantics for modal predicate logic.

\paragraph{What is proved.}
Nine claims are machine-checked.
The capstone definition
\code{mettaIL\_2modeDocFrag : TwoModeAdjointDocFrag}
(\codepath{MATTFragment.lean}) packages all of them into a single
kernel-checked Lean object; the individual components are listed in
Table~\ref{tab:matt}.

\begin{table}[h]
\centering
\begin{tabular}{lll}
\toprule
Claim & Lean declaration & Status \\
\midrule
Per-language modal Galois connection &
  \code{doctrine\_galois\_is\_langGalois} & proved \\
Per-language modal adjunction &
  \code{doctrine\_adjunction\_is\_langModalAdjunction} & proved \\
Eq-category pullback functoriality &
  \code{eqCategory\_mapPred\_functorial} & proved \\
Mode-skeleton composition/predicate laws &
  \code{mode2SkeletonLaws} & proved \\
Runtime/runtime commuting square &
  \code{runtime\_runtime\_square\_coherence} & proved \\
Runtime/behavioral commuting square &
  \code{runtime\_behavioral\_square\_coherence} & proved \\
Runtime morphism diamond witness transport &
  \code{runtime\_mode\_diamond\_transport\_comp} & proved \\
Pure mode isolation &
  \code{pure\_mode\_isolation} & proved \\
mettaPure runtime$\to$behavioral transport &
  \code{mettaPure\_runtime\_behavioral\_transport} & proved \\
\textbf{Capstone: MeTTa-IL is \code{TwoModeAdjointDocFrag}} &
  \code{mettaIL\_2modeDocFrag} & \textbf{proved} \\
\midrule
Full mode-2-category (2-cells/2-morphisms) & --- & out of scope \\
Syntactic MTT ($\square_\mu A$, locks, context ext.) & --- & out of scope \\
Pure-mode morphism theory & --- & deferred \\
\bottomrule
\end{tabular}
\caption{MaTT claim map (doctrinal/fibrational sense).  The proved-count is
  machine-checked: \code{provenCount\_eq : countByStatus .proven = 9}
  (\texttt{MATTClaimMap.lean}, closed by \code{decide}).}
\label{tab:matt}
\end{table}

\paragraph{The capstone theorem.}
A \code{TwoModeAdjointDocFrag} packages:
(i)~a mode category (identity, associativity laws for $\{$\texttt{pure},
\texttt{runtime}($L$), \texttt{behavioral}($L$)$\}$);
(ii)~a contravariant indexed predicate functor (\code{ModeHom.mapPred},
with identity and composition laws);
(iii)~a modal adjunction at each behavioral mode ($\lozenge \dashv \square$,
the OSLF Galois connection); and
(iv)~Beck-Chevalley commuting squares connecting the mode-level and
category-level predicate pullbacks.
The definition \code{mettaIL\_2modeDocFrag : TwoModeAdjointDocFrag}
fills all fields from existing proved lemmas.
The witness theorem \code{mettaIL\_is\_2mode\_adjoint\_doctrine :
Nonempty TwoModeAdjointDocFrag} is the Lean-kernel-checked statement
that this object exists.

\paragraph{Scope boundary.}
The out-of-scope items are explicit design choices.
Full 2-category structure requires 2-cells (natural transformations between
mode morphisms) and is deferred.
Syntactic MTT modal type formers require a full dependent kernel with
mode-indexed context extension; this is the role of MeTTa-Pure
(Appendix~\ref{app:metta-dtt}).
Pure-mode morphism theory is deferred pending the MeTTa-Pure bridge.

\paragraph{Reproducibility.}
\begin{lstlisting}[language={},caption={Build command for MaTT targets}]
ulimit -Sv 6291456 && LAKE_JOBS=3 nice -n 19 lake build \
  Mettapedia.OSLF.Framework.MATTFragment \
  Mettapedia.OSLF.Framework.MATTClaimMap \
  Mettapedia.OSLF.Framework.MATTProvableNow
\end{lstlisting}

The reviewer anchor in
\codepath{Mettapedia/OSLF/SpecIndex.lean}:265 is labelled
``Conservative MaTT Claim Map (No Overclaim)''.

\section{What This Companion Covers and Does Not Cover}

\textbf{Covered now:}
\begin{itemize}
\item MeTTaIL syntax + declarative/executable semantics bridge (sound and complete).
\item \code{languageDef!} macro with block-syntax parsing, carrier-aware types,
  authored binder preservation, syntax metasyntax operators, semantic validation,
  fail-fast core bridge, and faithful export.
\item RhoCalc and Metamath as acceptance-test \code{LanguageDef} instances
  (process fragment and core-step fragment, respectively).
\item MeTTaCore state machine + core progress/confluence-facing properties.
\item MeTTaMinimal/MeTTaFull OSLF instance endpoints.
\item HE MeTTa: computable recursive interpreter, 37 conformance theorems,
  OSLF \code{LanguageDef} + \code{PremiseProgram} export.
\item PeTTa: inductive big-step evaluation, LP soundness, grounded oracle,
  typed evaluation, effect tracking.
\item Governance refinement: shared \code{MeTTaLike} interface, HE/PeTTa
  DTS bridges, let* unfolding, canonical conformance (OB$\Rightarrow$PE).
\item Executable Rust backends: MeTTaHE (from Lean pipeline), MeTTaFullState
  (legacy), PeTTa (alias wrapper), all with REPL + oracle integration.
\item MeTTa-IL as 2-mode adjoint doctrine fragment (MaTT, 9 proved claims).
\end{itemize}

\textbf{Not yet full end-to-end by construction:}
\begin{itemize}
\item Full RhoCalc \code{languageDef!} parity: COMM and Extrude require
  premise-side \code{*map(...)} and pattern-side \code{*zip(...).*map(...)}.
\item Metamath \code{LanguageDef}-to-\code{mm-lean4} conformance theorem
  (type correspondence exists; semantic bridge is future work).
\item Fully automatic Lean-to-Rust export without any manual adaptation
  (field-name hygiene and builtin rendering still require hand-editing).
\item Standalone PeTTa Rust backend (currently delegates to MeTTaFullState).
\item Full HE test suite validation against the Rust MeTTaHE backend.
\item Full Hyperon atomspace/runtime integration (beyond encoded relation
  environment subset).
\end{itemize}

\section{Immediate Next Steps}

\begin{enumerate}
\item Extend \code{languageDef!} with premise-side \code{*map(...)} and
  pattern-side \code{*zip(...).*map(...)} to close the RhoCalc COMM/Extrude gap.
\item Prove the Metamath conformance bridge: \code{LanguageDef} rewrite semantics
  agree with \code{mm-lean4}'s \code{verify\_parser\_acceptance\_iff\_spec\_provable}.
\item Run the HE test suite against the Rust MeTTaHE backend to validate
  semantic correctness of the generated premise rules.
\item Build a standalone PeTTa Rust backend from the Lean PeTTa specs,
  replacing the current MeTTaFullState delegation.
\item Fix the \code{language!} proc macro to use hygienic variable names
  in generated parser code, preventing field-name collisions.
\item Tie this note into \code{leanOSLF.tex} as a companion reference once stable.
\end{enumerate}

%% ============================================================
\appendix
\section{WIP: Toward MeTTa-Pure --- A Dependent Kernel for MeTTa}
\label{app:metta-dtt}
%% ============================================================

\emph{This appendix sketches a forward-looking architecture for adding a
trusted dependent type theory (DTT) layer to MeTTa without changing the
semantics of existing MeTTa programs. It is work-in-progress design notes,
not a finalized specification.}

\subsection{Motivation: Two Results in This Paper}

Section~\ref{sec:sr} proves that MeTTa's current \code{HasType} judgment does
not satisfy subject reduction. This is not a bug; it is a precise consequence
of a design choice: \code{HasType} is annotation lookup, not inductive
inference. No single-point patch can fix this while preserving the rest of
MeTTa's open-world rewriting semantics.

The productive response is architectural, not remedial: \emph{build a small
dependent kernel alongside MeTTa, not inside it}.

\subsection{The Four-Layer Stack}

A clean MeTTa-DTT separates concerns across four layers:

\medskip
\begin{tabular}{lp{9cm}}
\toprule
Layer & Role \\
\midrule
\textbf{A — MeTTa-Pure} &
  A tiny intensional DTT kernel: ordered contexts $\Gamma$, $\Pi/\Sigma$-types,
  identity types, universes, inductive families. Bidirectionally checked.
  Subject reduction and decidable conversion hold by construction. \\[4pt]
\textbf{B — Elaboration boundary} &
  A (not necessarily trusted) elaborator translating MeTTa surface syntax into
  kernel terms. A checked reflection bridge maps quoted MeTTa atoms to kernel
  terms and back. \\[4pt]
\textbf{C — Full MeTTa runtime} &
  Current MeTTa unchanged: equations, atomspaces, \code{match},
  \code{superpose}, \code{collapse}, grounded interop, spaces,
  nondeterminism, concurrency. The AGI language proper. \\[4pt]
\textbf{D — OSLF overlay} &
  Behavioral/spatial/modal types derived from the operational semantics of
  Layer C, \emph{not} from kernel conversion. The modal adjunction
  $\Diamond \dashv \Box$ and the presheaf/topos lift remain exactly here. \\
\bottomrule
\end{tabular}
\medskip

The core thesis: \textbf{DTT for the trusted structural core; OSLF for
behavior; full MeTTa for open-world agent computation.}

\subsection{Why MeTTaIL Is the Right Substrate}

The kernel should be \emph{specified and implemented in MeTTaIL}, not beside it.
MeTTaIL already provides:
\begin{itemize}
  \item locally nameless binder representation with cofinite quantification;
  \item proved substitution lemmas (\code{subst\_open}, \code{subst\_lc}, etc.);
  \item declarative one-step relation \code{DeclReducesWithPremises} with
        soundness/completeness against the executable engine;
  \item OSLF \code{langOSLF} synthesis from any \code{LanguageDef}.
\end{itemize}
A MeTTa-Pure \code{LanguageDef} is therefore a \emph{language instance} of
the existing MeTTaIL framework, not a new implementation. The
\code{mettaPureOSLF} type system would be a fourth instance alongside
\code{mettaMinimalOSLF}, \code{mettaFullOSLF}, and \code{tinyMLOSLF}.

\subsection{Kernel Judgment Forms}

The kernel uses separate judgment forms distinguished from atomspace typing:

\begin{lstlisting}[language={},caption={Proposed MeTTa-Pure judgment forms}]
-- Well-formed type at universe level i
Sk ; G |- A  type  i

-- Synthesis: infer type of t
Sk ; G |- t  =>  A

-- Checking: verify t against A
Sk ; G |- t  <=  A

-- Definitional equality
Sk ; G |- t  ==  u  :  A
\end{lstlisting}

The local context \code{G} is an \emph{ordered list} of typed variables
\code{[(x1,A1), ..., (xn,An)]}---not an unordered atomspace.
The signature \code{Sk} carries declared kernel constants.

\subsection{Surface Syntax (S-Expression Style)}

The kernel surface can remain MeTTa-shaped:

\begin{lstlisting}[language={},caption={Proposed MeTTa-Pure surface forms}]
(Pi   (($x A)) B)       -- dependent function type
(lam  (($x A)) t)       -- lambda abstraction
(app  t u)              -- application
(Sg   (($x A)) B)       -- dependent pair type
(pair a b)              -- pair introduction
(fst p)  (snd p)        -- pair elimination
(Id   A a b)            -- identity type
(refl a)                -- reflexivity
(Type 0)  (Type 1) ...  -- universe hierarchy
\end{lstlisting}

This keeps the S-expression aesthetic while elaborating to a real binder AST.
At Layer C, \code{=} stays the runtime rewrite operator; the surface \code{:}
annotation gains a second meaning as a kernel declaration when inside a
pure context.

\subsection{Kernel Conversion: What Is and Is Not Allowed}

Kernel definitional equality is deliberately minimal:

\begin{itemize}
  \item $\beta$: $(\lambda x.\, t)\, u \equiv t[u/x]$ for $\Pi$/$\Sigma$ eliminations.
  \item $\delta$: transparent kernel constant unfolding.
  \item $\iota$: reduction rules for inductive eliminators.
  \item $\zeta$: local \code{let}-unfolding (optional, guarded by termination).
  \item $\eta$: function and pair extensionality, only where clearly desired.
\end{itemize}

\textbf{Not in kernel conversion:}
\begin{itemize}
  \item Arbitrary user-defined \code{=} rewrite rules (Layer C only).
  \item Atomspace lookup or mutation.
  \item \code{match}/\code{superpose}/\code{collapse} forms.
  \item Python/native grounded calls.
\end{itemize}

If rewrite rules are later admitted to definitional equality, this requires a
separate certified rewrite subsystem with a machine-checked
confluence/termination certificate---the $\lambda\Pi$-modulo route
(Dedukti/Saillard)---as a distinct upgrade, not a default.

\subsection{The ``Don'ts''}

Four constraints characterize the kernel boundary:
\begin{enumerate}
  \item \textbf{Do not} let arbitrary MeTTa rewrite rules define kernel conversion.
  \item \textbf{Do not} use an unordered atomspace as the kernel context.
  \item \textbf{Do not} identify OSLF native types with dependent types; the
        two layers are orthogonal.
  \item \textbf{Do not} allow full runtime nondeterministic evaluation inside
        types; type-level computation must terminate.
\end{enumerate}

\subsection{Categorical Summary}

The three-part categorical picture is:

\begin{center}
\begin{tabular}{lp{8cm}}
\toprule
Layer & Categorical home \\
\midrule
Kernel (MeTTa-Pure) &
  Categories with Families (CwF); contextual categories;
  locally cartesian closed categories; presheaf and groupoid models. \\[4pt]
Full MeTTa / OSLF &
  Presheaf/topos/modal/fibrational semantics over the operational rewrite
  system; the $\Diamond \dashv \Box$ modal adjunction. \\[4pt]
Bridge &
  Quotation and checked reflection between kernel terms and full MeTTa atoms. \\
\bottomrule
\end{tabular}
\end{center}

The presheaf topos semantics of OSLF (already in the Mettapedia formalization)
wraps a CwF kernel---exactly the structure that makes the categorical split
coherent rather than arbitrary.

\subsection{Implementation Reference Points}

The closest existing architectures are:
\begin{itemize}
  \item \textbf{Lean 4}: small total kernel, bidirectional elaborator, erasure
        for runtime performance; \code{partial} definitions are kernel-opaque.
  \item \textbf{Isabelle/Pure}: tiny LF-style framework underneath object
        logics (HOL, ZF, \ldots). Pure is itself a minimal DTT.
  \item \textbf{Isabelle/Spartan}: DTT as an object logic inside Pure,
        directly analogous to MeTTa-Pure as an object logic inside MeTTaIL.
\end{itemize}

In one sentence: \emph{a clean MeTTa-DTT is the smallest conservative
extension of MeTTaIL that adds a CwF-style dependent kernel with decidable
conversion and subject reduction, while treating today's MeTTa as an effectful
reflective metalanguage over that kernel, and OSLF as the behavioral logic of
the resulting operational world.}

\section{MQ-Calculus Lean Walkthrough}
\label{app:mq-lean}

This appendix maps paper-level MQ clauses to concrete Lean definitions.
All listings are excerpted verbatim from the formalization.

\paragraph{Files.}
\begin{itemize}
\item \codepath{MQCalculus/Syntax.lean} --- process grammar and gate types
\item \codepath{MQCalculus/CommRule.lean} --- Born-rule probability and COMM rule
\item \codepath{MQCalculus/Reduction.lean} --- full structural reduction relation
\item \codepath{MQCalculus/Backend.lean} --- abstract semantic backend interface
\item \codepath{MQCalculus/Denotational.lean} --- backend-parametric denotation
\item \codepath{MQCalculus/PaperMap.lean} --- stable audit surface (paper $\to$ theorem)
\end{itemize}

\subsection*{A. Process grammar and Born-rule reduction}

\begin{lstlisting}[language={},caption={Process grammar (\texttt{Syntax.lean}) and COMM rule (\texttt{CommRule.lean})}]
-- Syntax.lean
inductive Process : Type where
  | MQNil  : Process
  | MQPar  : Process -> Process -> Process
  | MQNu   : Process -> Process                    -- new-wire binder
  | MQGate : GateSpec -> Process -> Process        -- typed gate, then P
  | MQOut  : Nat -> Process                        -- emit wire i
  | MQIn   : Nat -> Process -> Process -> Process  -- measure i: 0 -> P, 1 -> Q

-- CommRule.lean
-- n-qubit Hilbert space as EuclideanSpace over Fin (2^n)
abbrev QVec (n : Nat) := EuclideanSpace Complex (Fin (2 ^ n))

-- Born-rule probability |<r|psi>|^2, proved (not axiomatized) from Mathlib
noncomputable def born_prob (r : Outcome) (psi : QVec 1) : Real :=
  match r with
  | .zero => Complex.normSq (inner ket0 psi)
  | .one  => Complex.normSq (inner ket1 psi)

-- Probabilities sum to 1 for unit-norm states (proved from PiLp.norm_sq_eq_of_L2)
theorem born_prob_sum_one (psi : QVec 1) (h : norm psi = 1) :
    born_prob .zero psi + born_prob .one psi = 1

-- COMM rule: MQOut i | MQIn i P Q can branch to either P or Q
inductive CommReduction : Nat -> Process -> Process -> MeasurementBranch -> Prop where
  | outcome_zero (i : Nat) (p q : Process) :
      CommReduction i p q { outcome := .zero, result := p }
  | outcome_one  (i : Nat) (p q : Process) :
      CommReduction i p q { outcome := .one,  result := q }
\end{lstlisting}

\subsection*{B. Reduction relation and backend-parametric denotation}

\begin{lstlisting}[language={},caption={Reduction (\texttt{Reduction.lean}) and denotation (\texttt{Denotational.lean})}]
-- Reduction.lean
-- Full structural reduction: COMM fires under SC, parallel, nu, gate contexts
inductive Reduces : Process -> Process -> Prop where
  | comm (i : Nat) (p q : Process) (b : MeasurementBranch) :
      CommReduction i p q b ->
      Reduces (MQPar (MQOut i) (MQIn i p q)) b.result
  | sc       (p p' q' q : Process) :
      SC p p' -> Reduces p' q' -> SC q' q -> Reduces p q
  | par_l    (p p' q   : Process) : Reduces p p' -> Reduces (MQPar p q)  (MQPar p' q)
  | par_r    (p q  q'  : Process) : Reduces q q' -> Reduces (MQPar p q)  (MQPar p q')
  | nu_step  (p p'     : Process) : Reduces p p' -> Reduces (MQNu p)     (MQNu p')
  | gate_step (g : GateSpec) (p p' : Process) :
      Reduces p p' -> Reduces (MQGate g p) (MQGate g p')

-- MQOut i is irreducible (cannot step)
theorem mq_out_irreducible (i : Nat) {q : Process} (h : Reduces (MQOut i) q) : False

-- Multi-step with notation p ->* q
inductive MultiStep : Process -> Process -> Prop where
  | refl (p : Process)   : MultiStep p p
  | step (p q r : Process) : Reduces p q -> MultiStep q r -> MultiStep p r

-- Denotational.lean
-- Backend-parametric denotation: same 6 clauses, any compliant backend
noncomputable def denoteWith (backend : MQSemanticsBackend) :
    Process -> QState -> Option QState
  | .MQNil,      st => some st
  | .MQOut _,    st => some st
  | .MQGate g p, st => denoteWith backend p (backend.applyGate g st)
  | .MQNu p,     st => denoteWith backend p (backend.allocFresh st)
  | .MQPar p q,  st => (denoteWith backend p st).bind (denoteWith backend q)
  | .MQIn i p q, st =>
      if backend.branchProb i .zero st >= backend.branchProb i .one st then
        denoteWith backend p (backend.collapse i .zero st)
      else
        denoteWith backend q (backend.collapse i .one st)
\end{lstlisting}

\subsection*{C. Paper-map audit surface}

\codepath{PaperMap.lean} exposes one named theorem per paper clause.
To audit a specific claim, import the file and check the corresponding theorem.

\begin{lstlisting}[language={},caption={Stable paper-to-theorem index (\texttt{PaperMap.lean})}]
-- Section 6.3 denotational clauses (all rfl proofs from simp lemmas)
theorem paper63_nil     : denoteWith backend MQNil      st = some st
theorem paper63_out     : denoteWith backend (MQOut i)  st = some st
theorem paper63_gate    : denoteWith backend (MQGate g p) st =
                            denoteWith backend p (backend.applyGate g st)
theorem paper63_new     : denoteWith backend (MQNu p)   st =
                            denoteWith backend p (backend.allocFresh st)
theorem paper63_par     : denoteWith backend (MQPar p q) st =
                            (denoteWith backend p st).bind (denoteWith backend q)
theorem paper63_in      : denoteWith backend (MQIn i p q) st =
                            if backend.branchProb i .zero st >= ... then ...

-- Norm-1 preservation under collapse (statevector backend, positive mass)
theorem paper63_statevector_collapse_norm_one_of_raw_pos
    (hraw : 0 < rawOutcomeProb i out st) :
    norm (statevectorBackend.collapse i out st).psi = 1

-- COMM denotation equals explicit branch-state selection
theorem paper63_comm_denote_eq_measurement_selection (i : Nat) (p q : Process) (st : QState) :
    denote (MQIn i p q) st = ...

-- Section 3 communication reduction facts
theorem paper3_comm_both_outcomes (i : Nat) (p q : Process) :
    CommReduction i p q {.zero, p} /\ CommReduction i p q {.one, q}

theorem paper3_comm_exhaustive (i : Nat) (p q : Process) (b : MeasurementBranch) :
    CommReduction i p q b -> b = {.zero, p} \/ b = {.one, q}

-- Cross-calculus coherence: MQ COMM branching <-> MORK binary fold branching
theorem paper_mork_mq_branching_equiv (i : Nat) (p q : Process) :
    (CommReduction i p q {.zero, p} /\ CommReduction i p q {.one, q}) <->
    (forall (fold : MORK.FoldStep) (_ : fold.isBinary), ...)
\end{lstlisting}

\end{document}
